
\begin{abstract}
\small
\emergencystretch=2em

Epilepsy is a neurological disorder that requires
specialized clinical assessment. Electroencephalography
(EEG) plays an important role in the evaluation of
abnormal electrical activity in the brain. However,
the interpretation of EEG recordings can be a
time-consuming process, particularly when dealing
with long-duration studies and large volumes of
signal data.

This thesis presents the design and implementation
of EpilexIA, a web-based platform developed to
support healthcare professionals in the management,
visualization, and analysis of electroencephalographic
studies. The platform incorporates an artificial
intelligence component designed to support the
identification and review of potential epileptiform
abnormalities.

The system was developed using React for the web
interface, Java and Spring Boot for backend services,
and PostgreSQL for data persistence. The artificial
intelligence subsystem uses Python-based components,
FastAPI for service integration, and scientific
libraries for EEG signal processing. The architecture
separates clinical information management from the
components responsible for automated signal analysis.

EpilexIA provides functionalities for patient and
EEG study management, visualization of signals in
EDF and CSV formats, signal filtering, sensitivity
adjustment, EEG montages, anomaly annotation, and
the preparation of medical reports. The platform
also incorporates mechanisms for managing user
access and sharing studies among authorized
healthcare professionals.

The evaluation included functional and integration
tests of the main clinical management and EEG
visualization components, as well as a contextual
assessment conducted by a medical specialist
at the Central Hospital of the Instituto de
Previsión Social (IPS). The results provided
evidence of the correct operation of the main
functionalities evaluated. However, the complete
execution and clinical performance of the artificial
intelligence subsystem could not be fully verified
within the scope of the tests performed.

In conclusion, EpilexIA provides a technological
platform to support the organization, visualization,
and analysis of EEG studies, establishing a foundation
for further validation and improvement of
artificial intelligence-assisted anomaly detection.

\end{abstract}
